%\abstracttitle
% Single spacing can be turned on for the abstract
%
% {\singlespacing
Optimization algorithms are routinely applied to problems where the
only available feedback is the quality of a proposed solution -- no
structural information, no gradients. Evaluating how well such
algorithms perform has traditionally relied on a single question: how
long until the best possible solution is found? This measure, while
theoretically convenient, says nothing about solution quality before
that moment arrives, and fails entirely when the budget runs out
before the optimum is reached.

This report addresses that gap. It develops a formal connection
between a trajectory-aware performance measure from reinforcement learning, 
cumulative regret, and an existing framework from evolutionary
computation that tracks how quickly algorithms improve across
different quality levels, dubbed ``Runtime Profiles''. The resulting identity shows that these two
perspectives, previously treated as separate, are mathematically
equivalent. This provides the ``combined anytime measure'' the
literature has called for but not yet delivered.

To validate and explore this connection, a purpose-built experimental
framework was developed, enabling reproducible, large-scale
comparisons of several algorithm families across a broad collection
of benchmark problems. These problems range from smooth,
well-understood landscapes to deceptive, hierarchical, and highly
rugged ones, each posing a different kind of challenge.

The experiments confirm the theoretical identity numerically and
reveal that regret trajectories expose meaningful differences between
algorithms even in cases where traditional measures are silent: when
no algorithm finds the optimum within the available budget, when
landscapes trap algorithms in suboptimal regions, or when two
algorithms arrive at the same endpoint via very different paths.
Across all settings, regret provides structured diagnostic
information that convergence-based summaries traditionally cannot.
